
\section{Generalized projective lines}

\begin{definition} \label{def:generalized projective line}
\uses{def:abstract complete curve}
A \emph{generalized projective line} is an  abstract complete curve $X$ satisfying the following conditions.
\begin{enumerate}
\item[(a)] Every divisor of degree $0$ is principal. In other words the map $\Pic(X) \to \Z$
is an isomorphism.
\item[(b)] We have $H^1(X, \cO_X) = 0$. 
\item[(c)] There exists a closed point $x \in X$ of degree $1$ with algebraically closed residue field.
\end{enumerate}
\end{definition}

In the following, let $X$ be a generalized projective line.

\begin{lemma} \label{L:section from complete flag}
\uses{def:generalized projective line,def:complete flag}
Let $F$ be a vector bundle on $X$ admitting a complete flag with slope sequence
$-1,0,\dots,0,1$ (with $n$ zeroes for some nonnegative integer $n$).
Then $H^0(X, F) \neq 0$.
\end{lemma}
\begin{proof}
\uses{cor:function field of curve}
Let $E$ be the ring $H^0(X, \cO_X)$, which by Corollary~\ref{cor:function field of curve} is a curve.
By condition (c) of Definition~\ref{def:generalized projective line}, we can choose a closed point $x \in X$ of degree $1$;
let $\kappa_x$ be the residue field of $x$.
By condition (a) of Definition~\ref{def:generalized projective line}, the ideal sheaf $J_x$ of $x$ is isomorphic to $\cO_X(-1)$.
We can thus construct the pushout $H$ as in the commutative diagram
\[
\begin{tikzcd}
0 \arrow{r} & J_x \arrow{r} \arrow{d} & F \arrow{r}\arrow{d} & F/J_x \arrow{r} \arrow[d,equal] & 0 \\
0 \arrow{r} & \cO_X \arrow{r} & H \arrow{r} & F/J_x \arrow{r} & 0
\end{tikzcd}
\]
Since $\Ext^1_X(\cO_X, \cO_X) = H^1(X, \cO_X)$ vanishes by condition (b) of Definition~\ref{def:generalized projective line}, $H$ admits a filtration with successive quotients $\cO_X^{\oplus (n+1)}, \cO_X(1)$.

We may assume that $H^0(X, H(-1)) = 0$, as otherwise $H$ contains the subbundle $\cO_X(1)$ whose intersection with $F$ has degree $\geq 0$ and hence
by (b) is either $\cO_X$ or $\cO_X(1)$.
Then $H^0(X, H)$ injects into $H^0(X, H/H(-1))$, which we reinterpret as the fiber $H_x$ of $H$ at $x$.
Using the exact sequence
\[
0 \to H^0(X, \cO_X^{\oplus n+1}) \to H^0(X, H) \to H^0(X, \cO_X(1)) \to H^1(X, \cO_X^{\oplus n}) = 0
\]
(where the last equality comes from (b)),
we express the image of $H^0(X, H) \to H_x$ 
as $E e_1 + \cdots + E e_{n+1} + \kappa_x e_{n+2}$
for some basis $e_1,\dots,e_{n+2}$ of $H_x$ over $\kappa_x$.

Meanwhile, the image of $H^0(X, F) \subset H^0(X, H)$ in $H_x$ is the intersection of the image of $H^0(X, H)$ with the image of $F_x \to H_x$; the latter is a codimension 1 subspace of $H_x$ and so can be written as the kernel of a nonzero $\kappa_x$-linear functional $\lambda\colon H_x \to \kappa_x$. 
Define $a \in E$, $b \in \kappa_x$ by
\[
(a,b) := \begin{cases} (0, 1) & \lambda(e_{n+2}) = 0 \\ (1, -\lambda(e_1)/\lambda(e_{n+2})) & \lambda(e_{n+2}) \neq 0; \end{cases}
\]
then $ae_1 + be_{n+2}$ is the image in $H_x$
of a nonzero element of $H^0(X, F)$.
\end{proof}

\begin{lemma} \label{lem:minimal filtration}
\uses{def:generalized projective line,def:complete flag}
Let $F$ be a vector bundle of rank $n$ on $X$.
Then $F$ admits a complete flag whose slope sequence $m_1,\dots,m_n$ has the following properties.
\begin{enumerate}
\item[(a)]
For $i=1,\dots,n-1$, we have $m_{i+1} \leq m_{i}+1$.
\item[(b)]
There do not exist indices $1 \leq i < j \leq n$ such that for some $m \in \Z$,
\[
(m_i,\dots,m_j) = (m-1, m,\dots,m,m+1).
\]
\end{enumerate}
\end{lemma}
\begin{proof}
\uses{lem:subbundle slope bound,L:section from complete flag}
For each $i  \in \{1,\dots,n-1\}$, 
Lemma~\ref{lem:subbundle slope bound}
imposes a uniform upper bound on $m_1 + \cdots + m_i$,
and in all cases $m_1 +\cdots + m_n = \deg(F)$.
Consequently, we can choose a complete flag to minimize 
\[
\sum_{i=1}^n i m_i 
= (n+1) \deg(F) -\sum_{i=1}^{n-1} (m_1 + \cdots + m_i).
\]
We prove that this complete flag has the desired property.

Suppose by way of contradiction that (a) fails for some $i$. 
Then $F_{i+1}/F_{i-1}(m_i-1)$ admits a complete flag with degree sequence
$-1, m_{i+1}-m_i$; we then obtain a subbundle $(F_{i+1}/F_{i-1})(-m_i-1)$
admitting a complete flag with degree sequence $-1, 1$.
By Lemma~\ref{L:section from complete flag}, $F_{i+1}/F_{i-1}(-m_i-1)$ admits $\cO_X$ as a (not necessarily saturated) subbundle; we thus obtain a new complete flag in which the terms
$m_i, m_{i+1}$ in the degree sequence are replaced by new terms $m_i', m_{i+1}'$ with $m_i' > m_i$ and $m_i' + m_{i+1}' = m_i + m_{i+1}$, contradicting minimality.

Suppose by way of contradiction that (b) fails for some $i,j$.
Then $(F_j/F_i)(-m)$ admits a complete flag with degree sequence $-1,0,\dots,0,1$.
By Lemma~\ref{L:section from complete flag} again, 
$F_j/F_i(-m)$ admits $\cO_X$ as a (not necessarily saturated) subbundle;
we thus obtain a new complete flag in which the terms of $m_i, \dots, m_j$ in the degree sequence
are replaced by new terms $m_i', \dots, m_j'$ with $m_i' > m_i$
and $m_i' + \cdots + m_k' \geq m_i + \cdots + m_k$ for $k=i+1,\dots,j$,
again contradicting minimality.
\end{proof}

\begin{lemma} \label{lem:sequence argument1}
\leanok
Let $n$ be a positive integer.
Let $m_1,\dots,m_n$ be a sequence of integers satisfying 
the following conditions.
\begin{enumerate}
\item[(a)]
For $i=1,\dots,n-1$, we have $m_{i+1} \leq m_{i}+1$.
\item[(b)]
There do not exist indices $1 \leq i < j \leq n$ such that for some $m \in \Z$,
\[
(m_i,\dots,m_j) = (m-1, m,\dots,m,m+1).
\]
\end{enumerate}
Then for $1 \leq i < j \leq n$,
\[
m_j \leq m_i+1 \qquad (1 \leq i < j \leq n).
\]
\end{lemma}
\begin{proof}
\leanok
We prove the claim by induction on $j-i$.
The base case $j-i=1$ is covered by (a).
For the induction step, suppose by way of contradiction that $m_j> m_i+1$ for some $i,j$.
Set $m := m_i + 1$.
For $k = i+1, \dots, j-1$ we have $m_{k} \leq m_i+1$ and $m_j \leq m_{k}+1$ by the induction hypothesis;
combining with $m_j \geq m_i + 2$ shows that both of these must be equalities
and so $m_k = m, m_j = m+1$.
Now $(m_i, \dots, m_j) = (m-1, m, \dots, m, m+1)$, which contradicts (b).
\end{proof}


\begin{lemma} \label{lem:sequence argument2}
\leanok
Let $n$ be a positive integer.
Let $m_1,\dots,m_n$ be a sequence of integers satisfying 
the following conditions:
\begin{align}
m_1 +\cdots+m_n = 0 & \label{eq:bound2} \\
m_1 + \cdots + m_i \leq 0 & \quad (i=1,\dots,n-1)  \label{eq:bound1}\\
m_j \leq m_i+1 & \quad (1 \leq i < j \leq n). \label{eq:bound3}
\end{align}
Then $m_1 = \cdots = m_n = 0$.
\end{lemma}
\begin{proof}
\leanok
We proceed by induction on $n$, with trivial base case $n=1$.
For $n > 1$, note that \eqref{eq:bound1} and \eqref{eq:bound2} together imply $m_n \geq 0$.
On the other hand, \eqref{eq:bound1} implies $\min_{i<n} \{m_i\} \leq 0$, and then \eqref{eq:bound3}
implies $m_n \leq 1$. Finally, if $m_n = 1$, then \eqref{eq:bound3} implies that $m_i \geq 0$ for $i=1,\dots,{n-1}$, but then $m_1 + \cdots + m_n \geq 1$ in violation of \eqref{eq:bound2}.
We deduce that $m_n = 0$; now the induction hypothesis applies to $m_1,\dots,m_{n-1}$ to finish.
\end{proof}

\begin{proposition} \label{prop:semistable slope 0}
\uses{def:generalized projective line,def:semistable bundle}
Every semistable bundle of degree $0$ on $X$ is trivial.
\end{proposition}
\begin{proof}
\uses{lem:minimal filtration,lem:sequence argument1,lem:sequence argument2}
Let $F$ be a semistable vector bundle of degree $0$ and rank $n$.
Choose a complete flag as in Lemma~\ref{lem:minimal filtration}.
By Lemma~\ref{lem:sequence argument1} and 
Lemma~\ref{lem:sequence argument2}, the degree sequence must be identically $0$.
Since $\Ext^1_X(\cO_X, \cO_X) = H^1(X, \cO_X)$ vanishes by condition (b) of Definition~\ref{def:generalized projective line}, it follows that $F$ is trivial.
\end{proof}

\begin{lemma} \label{lem:one-step modification induction}
\uses{def:generalized projective line,def:HN filtration}
Let $F$ be a vector bundle on $X$ of rank $n$ with all HN slopes in $[0,1)$.
Let $x \in X$ be a closed point of degree $1$
and let $j_x \colon x \to X$ the canonical inclusion.
Then for $i=0,\dots,d$,
there exist a coherent sheaf $G_i$ on $x$ and an exact sequence
\begin{equation} \label{eq:modification}
0 \to H_i \to F \to j_x^* G_i \to 0
\end{equation}
of coherent sheaves on $X$ in which $\deg(H_i) = d-i$ and the HN slopes of $H_i$ are all nonnegative.
\end{lemma}
\begin{proof}
Each $i$-dimensional quotient $G_i$ of the fiber $F_x$ gives rise to an exact sequence as in \eqref{eq:modification} with $\deg(H_i) = d_i$. It thus suffices to enforce that $H_i$ has all HN slopesn nonnegative.
As the base case $i=0$ is vacuously true, we proceed to the induction step. To construct $G_i$,
let $K$ be the maximal submodule of $H_{i-1}$ of slope $0$; more precisely, $K$ is the first step of the HN filtration of $H_{i-1}$ if 0 occurs as an HN slope of $H_{i-1}$ and the zero subbundle otherwise. It then suffices to choose $G_i$ so that $K \subseteq H_i$, which is possible because $K \neq H_{i-1}$ (as otherwise $\deg(H_{i-1})=d-i+1$ would be zero).
\end{proof}

\begin{lemma} \label{lem:one-step modification}
\uses{def:generalized projective line,def:HN filtration}
Let $F$ be a vector bundle on $X$ with all HN slopes in $[0,1)$.
Let $x \in X$ be a closed point of degree $1$.
Then for $j_x \colon x \to X$ the canonical inclusion,
there exist a coherent sheaf $G$ on $x$ and an exact sequence
\[
0 \to \cO_X^{\oplus \rank(F)} \to F \to j_x^* G \to 0
\]
of coherent sheaves on $X$.
\end{lemma}
\begin{proof}
\uses{prop:semistable slope 0,lem:one-step modification induction}
Apply Lemma~\ref{lem:one-step modification induction} with $i = d :=\deg(F)$ to construct an exact sequence as in \eqref{eq:modification} in which $\deg(H_d) = 0$ and the HN slopes of $H_d$ are all nonnegative. Then $H_d$ is semistable of degree $0$, and hence trivial by Proposition~\ref{prop:semistable slope 0}.
\end{proof}

\begin{corollary} \label{cor:slopes to sections}
\uses{def:generalized projective line,def:HN filtration}
Let $F$ be a nonzero vector bundle on $X$ whose HN slopes are all nonnegative.
Then $F$ is generated by global sections and $H^1(X, F) = 0$.
\end{corollary}
\begin{proof}
\uses{lem:one-step modification}
By replacing $F$ with the first step of its HN filtration, we may reduce to the case where $F$ is semistable. By twisting, we may further ensure that $\mu(F) \in [0,1)$. We then deduce both claims from Lemma~\ref{lem:one-step modification} and the assumption $H^1(X, \cO_X) = 0$ from Definition~\ref{def:generalized projective line}.
\end{proof}

\begin{lemma} \label{lem:end of semistable}
\uses{def:generalized projective line,def:semistable bundle}
For any semistable vector bundle $F$ on $X$, $F^\vee \otimes F$ is trivial.
\end{lemma}
\begin{proof}
\uses{cor:slopes to sections,lem:map of stable,prop:semistable slope 0}
Using the assumption that $H^1(X, \cO_X) = 0$,
we reduce to the case where $F$ is stable with slope in $[0,1)$. In this case, suppose by way of contradiction that $F^\vee \otimes F$ is not semistable, and let $G$ be the first step of its HN filtration. Since $G \otimes G$ has slope $2\mu(G) > \mu(G)$, 
the composition map $(F^\vee \otimes F) \otimes (F^\vee \otimes F) \to F^\vee \otimes F$
restricts to zero on $G \otimes G$.
By Corollary~\ref{cor:slopes to sections}, $G$ is generated by global sections;
consequently, $G \otimes G$ is generated by nonzero global sections of the form $s_1 \otimes s_2$.
Any such section corresponds to a pair of nonzero endomorphisms $F \to F$ which compose to zero; this contradicts Lemma~\ref{lem:map of stable}.
We deduce that $F^\vee \otimes F$ is semistable of slope $0$,
and hence trivial by Proposition~\ref{prop:semistable slope 0}.
\end{proof}

\begin{proposition} \label{prop:tensor product of semistable}
\uses{def:generalized projective line,def:semistable bundle}
The tensor product of any two semistable vector bundles on $X$ is semistable.
\end{proposition}
\begin{proof}
\uses{lem:end of semistable,lem:semistable from product}
Let $F,G$ be two semistable vector bundles on $X$.
By Lemma~\ref{lem:end of semistable}, 
\[
(F^\vee \otimes G^\vee) \otimes (F \otimes G)
\cong (F^\vee \otimes F) \otimes (G^\vee \otimes G)
\]
is semistable. By Lemma~\ref{lem:semistable from product}, $F \otimes G$ must be semistable.
\end{proof}

\begin{lemma} \label{cor:unique stable}
\uses{def:generalized projective line,def:stable bundle}
For any $\lambda \in \Q$, there is at most one isomorphism class of stable vector bundles on $X$ of slope $\lambda$.
\end{lemma}
\begin{proof}
\uses{prop:tensor product of semistable,lem:map of stable}
Let $F,G$ be two such bundles. By Proposition~\ref{prop:tensor product of semistable}, $F^\vee \otimes G$ is semistable; since it has slope 0, it is trivial by Proposition~\ref{prop:semistable slope 0}.
We can thus find a nonzero morphism $f\colon 
F \to G$, which must be an isomorphism by Lemma~\ref{lem:map of stable}.
\end{proof}

\begin{lemma} \label{cor:split sequence in order}
\uses{def:generalized projective line,def:HN filtration}
Any exact sequence of vector bundles on $X$ of the form
\[
0 \to F_1 \to F \to F_2 \to 0,
\]
in which every HN slope of $F_1$ is greater than or equal to every HN slope of $F_2$, is split.
\end{lemma}
\begin{proof}
\label{uses:prop:tensor product of semistable,cor:slopes to sections}
By Proposition~\ref{prop:tensor product of semistable}, $F_2^\vee \otimes F_1$ has nonnegative HN slopes.
Now $\Ext^1_X(F_2, F_1) = H^1(X, F_2^\vee \otimes F_1)$ which vanishes by Corollary~\ref{cor:slopes to sections}.
\end{proof}

\begin{proposition} \label{prop:decompose as stables}
\uses{def:generalized projective line,def:HN filtration,def:stable bundle,def:polystable bundle}
Every vector bundle on $X$ splits (nonuniquely) as a direct sum of stable bundles. In particular, the HN filtration always splits and every semistable bundle is polystable.
\end{proposition}
\begin{proof}
\uses{cor:split sequence in order}
Both assertions follow directly from Lemma~\ref{cor:split sequence in order}.
\end{proof}

\begin{lemma} \label{lem:construct stable degree 1}
\uses{def:generalized projective line,def:stable bundle,cor:function field of curve}
Let $E$ be the field $H^0(X, \cO_X)$ (Corollary~\ref{cor:function field of curve}).
For $d$ a positive integer with $\dim_{E} H^0(X, \cO_X(1)) > d$,
there exists a semistable vector bundle on $X$ of rank $d$ and degree $1$.
\end{lemma}
\begin{proof}
Let $x$ be a closed point of $X$ of degree 1.
Set $F := \cO_X^{\oplus d}$ and let $e_1,\dots,e_d$ be the images in $F_x$
of the standard basis vectors of $H^0(X, F) \cong E^{d}$.
The condition $\dim_{E} H^0(X, \cO_X(1)) > d$ is equivalent to $\dim_{E} \kappa_x \geq d$;
consequently, we can choose a one-dimensional quotient $V$ of $F_x$
for which the composition $E^d \to F_x \to V$ is injective.
Construct the subbundle $G$ of $F$
so that we have an exact sequence
\[
0 \to G \to F \to j_x^* V \to 0.
\]
We then have another exact sequence
\[
0 \to F \to G^\vee \to j_x^* V^\vee \to 0
\]
using which we see that $G^\vee$ is stable of degree $1$.
\end{proof}

\begin{lemma} \label{lem:construct semistables}
\uses{def:generalized projective line,def:semistable bundle,cor:function field of curve}
Let $E$ be the field $H^0(X, \cO_X)$ (Corollary~\ref{cor:function field of curve}).
Suppose that $\dim_{E} H^0(X, \cO_X(1)) = \infty$.
Then for every positive integer $d$
and every $c \in \Z$ coprime to $d$, there exists a stable vector bundle on $X$ of rank $d$ and degree $c$.
\end{lemma}
\begin{proof}
\uses{prop:decompose as stables,lem:construct stable degree 1,cor:unique stable}
We proceed by induction on $d$, the case $d=1$ being evident because $\cO_X(c)$ is stable for all $c \in \Z$. 

We will produce several examples of semistable bundles $F$ of slope $\frac{c}{d}$. For each such bundle, 
we may apply Proposition~\ref{prop:decompose as stables} to decompose $F$ as a sum of stable bundles.
Let $G$ be one such bundle; by Lemma~\ref{cor:unique stable}, all of the other summands are isomorphic to $G$, and so $\rank(G)$ divides $\rank(F)$. By Lemma~\ref{cor:unique stable} again, the isomorphism class of $G$ does not depend on the choice of $F$ either. It will thus suffice to produce enough examples so that the gcd of the resulting values of $\rank(F)$ is equal to $d$.

\begin{itemize}
\item
For $d>1$, we can choose $b \in \{1,\dots,d-1\}$ with $bc \equiv 1 \pmod{d}$.
By Lemma~\ref{lem:construct stable degree 1}, we can find a stable bundle $G_1$ of rank $bd$ and degree $1$. By the induction hypothesis, we can find a stable bundle $G_2$ of rank $b$ and degree $(bc-1)/d$. 
By Proposition~\ref{prop:tensor product of semistable}, $G_1 \otimes G_2$ is semistable of rank $b^2 d$ and slope
\[
\frac{1}{bd} + \frac{(bc-1)/d}{b} = \frac{c}{d}.
\]

\item
By Lemma~\ref{lem:construct stable degree 1} again, we can find a stable bundle $G_1'$ of rank $(d-b)d$ and degree $1$. By the induction hypothesis, we can find a stable bundle $G_2'$ of rank $d-b$ and degree $((d-b) c+1)/d$. 
By Proposition~\ref{prop:tensor product of semistable}, $(G_1')^\vee \otimes G_2$ is semistable of rank $(d-b)^2 d$ and slope
\[
-\frac{1}{(d-b)d} + \frac{((d-b)c+1)/d}{d-b} = \frac{c}{d}.
\]
By Proposition~\ref{prop:decompose as stables}, $(G_1')^\vee \otimes G_2$ decomposes as a sum of stable bundles  of slope $\frac{c}{d}$. Let $G'$ be one such bundle; by Lemma~\ref{cor:unique stable}, all of the other summands are isomorphic to $G'$. We deduce that $\rank(G')$ divides $(d-b)^2 d$.

\item
By Lemma~\ref{lem:construct stable degree 1} again, we can find a stable bundle $G''_1$ of rank $d$ and degree $1$. 
By Proposition~\ref{prop:tensor product of semistable}, $(G''_1)^{\otimes c}$ is semistable of rank $d^c$ and slope $\frac{c}{d}$.
\end{itemize}
We deduce that there exists a stable bundle $G$ of slope $\frac{c}{d}$ and degree dividing
\[
\gcd(b^2 d, (d-b)^2 d, d^c) = d \gcd(b^2, (d-b)^2, d^{c-1}) = d.
\]
On the other hand, since $\deg(G) \in \Z$, $\rank(G)$ must be divisible by $d$.
Hence $G$ is the desired stable bundle of rank $d$ and degree $c$.
\end{proof}

 
\begin{lemma} \label{lem:apply artin-schreier}
\uses{def:generalized projective line,cor:function field of curve}
Let $E$ be the field $H^0(X, \cO_X)$ (Corollary~\ref{cor:function field of curve}).
Then one of the following conditions hold.
\begin{enumerate}
\item[(a)]
We have $\dim_E H^0(X, \cO_X(1)) = 2$ and $X \cong \mathbf{P}^1_E =\Proj E[x,y]$ (the \emph{classical case}). In this case, $H^1(X, \cO_X(-1)) = 0$.
\item[(b)]
We have $\dim_E H^0(X, \cO_X(1)) = 3$ and $X \cong \Proj E[x,y,z]/(x^2+y^2+z^2)$ (the \emph{twistor case}). In this case, there is an isomorphism
$X \times_E E(\sqrt{-1}) \cong \mathbf{P}^1_{E(\sqrt{-1})}$ via which $\cO_X(1)$ corresponds to $\cO_{\mathbf{P}^1}(2)$.
\item[(c)]
We have $\dim_E H^0(X, \cO_X(1)) = \infty$.
\end{enumerate}
\end{lemma}
\begin{proof}
\uses{thm:artin schreier thm}
Apply Theorem~\ref{thm:artin schreier thm}.
\end{proof}

\begin{theorem} \label{thm:abstract classification}
\uses{def:generalized projective line,def:stable bundle}
Set
\[
S := \begin{cases} \Z & \mbox{in the classical case,} \\
\tfrac{1}{2} \Z & \mbox{in the twistor case,} \\
\Q & \mbox{otherwise}.
\end{cases}
\]
Then the following statements hold.
\begin{enumerate}
\item[(a)]
There exists exactly one isomorphism class of stable vector bundles $\cO_X(\lambda)$ of slope $\lambda$
if $\lambda \in S$ and none if $\lambda \notin S$.
\item[(b)]
For $\lambda \in S$ with denominator $d$ in lower terms, $\rank(\cO_X(\lambda)) = d$.
\item[(c)]
The endomorphisms of $\cO_X(\lambda)$ form a division algebra of degree $d$ over $P_0$.
\item[(d)]
For $\lambda, \lambda' \in S$, $\cO_X(\lambda) \otimes \cO_X(\lambda') \cong \cO_X(\lambda+\lambda')^{\oplus m}$
for some positive integer $m$.
\item[(d)]
Every vector bundle on $X$ splits (nonuniquely) as a direct sum of summands each of the form $\cO_X(\lambda)$ for some $\lambda \in S$.
\end{enumerate}
\end{theorem}
\begin{proof}
\uses{lem:construct semistables,cor:unique stable,lem:apply artin-schreier,lem:filtration by line bundles,lem:end of semistable,prop:tensor product of semistable,prop:decompose as stables}
If $\dim_{P_0} P_1 = \infty$, then 
we may deduce (a) and (b) from Lemma~\ref{lem:construct semistables} and Lemma~\ref{cor:unique stable}. Otherwise, by Lemma~\ref{lem:apply artin-schreier} we may argue as follows.
\begin{itemize}
\item 
In the classical case, $H^1(X, \cO_X(-1)) = 0$.
Hence any exact sequence of the form
\[
0 \to \cO_X \to F \to \cO_X(1) \to 0 
\]
splits, as then does any filtration as in Lemma~\ref{lem:minimal filtration}, proving (a).
From this, (b) is obvious.
\item
In the twistor case, by identi
by identifying $P \otimes_{P_0} P_0(\sqrt{-1})$ with $P_0(\sqrt{-1})[x^2,xy,y^2]$ and applying the classical case, we see that $S \subseteq \tfrac{1}{2}\Z$.
We deduce (a) and (b) by applying Lemma~\ref{lem:construct stable degree 1} to construct a stable bundle of rank $2$ and degree $1$, then applying Lemma~\ref{cor:unique stable} for uniqueness.
\end{itemize}

Given (a)--(b), (c) follows from Lemma~\ref{lem:map of stable} and Lemma~\ref{lem:end of semistable}; 
(d) is immediate from Proposition~\ref{prop:tensor product of semistable} and Proposition~\ref{prop:decompose as stables};
and (e) is immediate from Proposition~\ref{prop:decompose as stables}.
\end{proof}
